\subsection{Ausblick: Nächste Schritte}

Der Prototypentest bestätigt die Kernidee des \enquote{ZVV Tageserlebnis} und zeigt zugleich, wo das Konzept vor einer breiten Einführung nachgeschärft werden muss. Der Ausblick bündelt diese offenen Punkte zu einem Weg, der über den unmittelbaren Pilotbetrieb hinausreicht und das Angebot schrittweise vom Prototyp zum festen Bestandteil des Freizeitangebots des \gls{zvv} entwickelt.

\subsubsection{Vom Prototyp zum Pilot}
Der nächste konkrete Schritt ist ein realer Pilotbetrieb mit den drei ausgearbeiteten Touren in einer ausgewählten Region. Erst der Feldtest mit echten Gastgebern und einer grösseren, breiteren Gruppe macht aus den qualitativen Tendenzen belastbare Erkenntnisse. Im Zentrum stehen dabei die beiden im Test sichtbar gewordenen Schwachstellen: die \textbf{Vorab-Bezahlung} muss vollständig digital und vor dem Ausflugstag gelöst sein, damit der Gastgeber am Treffpunkt nicht mit dem Geldeinzug beschäftigt ist, und der \textbf{Zugang} muss niederschwellig bleiben, mit einfachem Weblink statt QR-Code und stets auch einem analogen Weg für Menschen ohne Smartphone.

\subsubsection{Offene Lücken im Erlebnis schliessen}
Bereits in der Develop-Phase wird sichtbar, dass das Konzept einzelne Etappen der Reise noch nicht vollständig abdeckt. Diese Lücken bleiben für die Weiterentwicklung zentral: der \textbf{Erstzugang} zu zurückgezogenen und vulnerablen Senioren, gerade im ländlichen Raum, die rein digital nicht erreichbar sind, die \textbf{erste und letzte Meile} von zuhause zum Treffpunkt und zurück, sowie die \textbf{Begleitung am Tagesende}, wenn die Energie nachlässt. Hier liegt das grösste Potenzial, das Angebot über den heutigen Prototyp hinaus zu vervollständigen.

\subsubsection{Skalierung und Rolle des ZVV}
Trägt der Pilot, lässt sich das Angebot in drei Stufen ausbauen: die \textbf{Integration des \gls{oev}-Tickets} in Tarif und Buchung, der \textbf{Ausbau der Touren} um weitere Themen, Bausteine und Regionen sowie der schrittweise \textbf{Rollout im Kanton} mit vielen Gastgebern. Entscheidend ist dabei die Rolle des \gls{zvv} als Plattform: Er liefert das kuratierte Erlebnis, die Tarif-Anbindung und die Gewinnung neuer Gastgeber, während die Botschafter wie Kurt die Begeisterung in ihr Umfeld tragen. So füllt das Angebot freie Sitze in den \gls{nvz}, ohne dass zusätzliche Infrastruktur nötig wird.

\subsubsection{Wirkung messbar machen}
Damit der Pilot mehr ist als ein Stimmungsbild, braucht es von Beginn an klare Kriterien. Sinnvoll messbar sind etwa die Zahl aktiver Gastgeber, die durchschnittliche Gruppengrösse, der Anteil der Teilnehmenden, die zuvor selten den \gls{oev} für Freizeitwege nutzten, sowie die Wiederbuchungsrate. Diese Grössen zeigen, ob das Angebot tatsächlich auf die strategischen Ziele der Design Challenge einzahlt, mehr Auslastung in den \gls{nvz}, eine Verschiebung der \gls{modalwahl} und langfristige Verhaltensänderung.

\subsection{Reflexion und Erkenntnisse der Arbeitsgruppe zum Prozess}

Über die einzelnen Phasen hinaus bringt uns die Projektarbeit als Gruppe wertvolle Erkenntnisse zum Prozess selbst. Die folgende Reflexion fasst zusammen, was im Vorgehen entlang des \gls{doublediam} besonders trägt, wo wir an Grenzen stossen und was wir für künftige Innovationsprojekte mitnehmen.

\subsubsection{Der Wert des doppelten Diamanten}
Die grösste Lernerfahrung ist die Disziplin, den Lösungsraum erst spät zu öffnen. Gerade die Versuchung, früh in vermeintlich fertige Lösungen zu springen, ist gross. Das bewusste Verweilen im Problemraum während der Discover- und Define-Phase zahlt sich aus: Erst das tiefe Verständnis der Zielgruppe macht den entscheidenden Perspektivwechsel möglich. Der Wechsel zwischen divergentem Öffnen und konvergentem Verdichten gibt dem Prozess dabei eine klare Struktur und hilft, in jeder Phase die richtige Frage zu stellen \parencite{designcouncil2005}.

\subsubsection{Das Reframing als Wendepunkt}
Die wichtigste inhaltliche Erkenntnis ist der Perspektivwechsel von der Define-Phase: Nicht die bereits überzeugte vielreisende Person ist der Hebel, sondern die Persona Kurt als \textbf{Botschafter} für ihr überredbares Umfeld. Diese Umdeutung über den \gls{pov} und die daraus abgeleiteten \gls{hmw}n trägt das gesamte spätere Konzept. Für uns ist es die zentrale Bestätigung, dass ein gut gewähltes Reframing mehr bewirkt als jede einzelne Lösungsidee. Auch die \gls{kopfstand} erweist sich als wirkungsvoll, weil die bewusste Umkehrung der Fragestellung Barrieren schneller sichtbar macht als die direkte Suche nach Lösungen.

\subsubsection{Testen heisst, sich überraschen lassen}
Die Deliver-Phase lehrt uns den Wert der Haltung, einen Prototyp zu bauen, als hätte man recht, ihn aber zu testen, als läge man falsch \parencite{leifer2015}. Erst die Konfrontation mit echten Nutzenden zeigt, dass nicht die aufwendig gestalteten Apps, sondern der einfache physische Treffpunkt am stärksten trägt, und dass die digitale Bezahlung eine ernsthafte Hürde ist. Beide Befunde widersprechen unserer ursprünglichen, app-zentrierten Erwartung, und genau darin liegt ihr Wert.

\subsubsection{Zusammenarbeit und Rahmenbedingungen}
Die Arbeit gelingt als gleichberechtigte Gruppenarbeit über alle Phasen hinweg, der gruppenübergreifende Austausch von Zwischenergebnissen und Interviews erweitert die Faktenbasis zusätzlich. Als produktiv erweist sich auch der enge Rahmen des \gls{zvv}, der Eingriffe in Pricing, Tarife und Fahrplan ausschliesst. Statt uns einzuschränken, lenkt diese Vorgabe die Kreativität auf die Hebel Kommunikation, Service, Erlebnis und persönliche Begleitung, jene Felder, in denen ein Verkehrsverbund tatsächlich Gestaltungsspielraum hat.

\subsubsection{Grenzen und was wir anders machen würden}
Selbstkritisch nehmen wir mit, dass die Validierung auf einer kleinen, qualitativen Stichprobe beruht und die Auswahl der Testpersonen auf einer plausiblen, aber nicht abschliessend geprüften Annahme fusst. In einem nächsten Projekt würden wir früher und mit einer breiteren Gruppe testen sowie einzelne Annahmen, etwa zur Zahlungsbereitschaft, gezielter im Feld prüfen, statt sie bis in die Deliver-Phase mitzutragen. Insgesamt zeigt uns die Arbeit, dass methodische Konsequenz, ein mutiges Reframing und das ehrliche Testen am Menschen gemeinsam zu einem Konzept führen, das eine echte Chance hat, vitale Senioren für den \gls{oev} zu begeistern.
